\documentclass[11pt]{article}
\usepackage[margin=1in]{geometry}
\usepackage{amsmath,amssymb,amsthm}
\usepackage[colorlinks=true,linkcolor=blue,citecolor=blue,urlcolor=blue]{hyperref}
\newtheorem{theorem}{Theorem}[section]
\newtheorem{lemma}[theorem]{Lemma}
\newtheorem{proposition}[theorem]{Proposition}
\newcommand{\R}{\mathbb R}
\newcommand{\rank}{\operatorname{rank}}
\newcommand{\tr}{\operatorname{tr}}
\newcommand{\diag}{\operatorname{diag}}
\newcommand{\spanof}{\operatorname{span}}
\newcommand{\norm}[1]{\left\lVert#1\right\rVert}
\title{An exact spectral obstruction for one edge change\protect\\
on a weighted path with unit vertex masses}
\author{Research working note}
\date{November 15, 2026}
\begin{document}
\maketitle

\begin{abstract}
A positive rank-one update can produce a fractional relative spectrum
that no single edge increase on a finite weighted path with unit vertex
masses realizes exactly, even after padding with zero correction
eigenvalues. We give the explicit square-root example
$(q_1,q_2)=(5/4,4\sqrt2/5)$ at leverage one. The proof classifies every
path whose correction has rank two and computes the exact minimum of
the sum of its two nontrivial relative eigenvalues. At leverage one
this minimum is $2.4080334082\ldots$, whereas the proposed target has
sum $2.3813708498\ldots$. An explicit two-dimensional positive definite
matrix pair attains the target. Approximation by longer paths with
additional small positive correction eigenvalues remains possible;
the obstruction concerns exact realization. Historical novelty is
unverified.
\end{abstract}

\section{Relative spectra and the obstruction}

For $A\succ0$, $B=A+uu^T$, and $0<\alpha<1$, put
\[
 Q_\alpha(A,u)=A^{-\alpha/2}B^\alpha A^{-\alpha/2},
 \qquad C_\alpha(A,u)=Q_\alpha(A,u)-I.
\]
The leverage is $\tau=u^TA^{-1}u$. All powers are principal powers.
Operator monotonicity and the determinant lemma give the classical
necessary conditions
\begin{equation}\label{eq:budget}
 C_\alpha\succeq0,\qquad
 \det Q_\alpha=(1+\tau)^\alpha.
\end{equation}

Let $L$ be the combinatorial Laplacian of a finite connected weighted
path. Every vertex has mass one, and every path-edge conductance is
strictly positive. Increasing an edge of conductance $w$ by $\tau w$
gives
\[
 L'=L+\tau wbb^T,\qquad b=e_j-e_{j+1}.
\]
Both matrices have kernel $\spanof\{\mathbf1\}$. Restricting to
$S=\mathbf1^\perp$ makes the baseline positive definite. Equivalently,
define the full relative correction by Moore--Penrose functional
calculus:
\begin{equation}\label{eq:graphcorrection}
 \mathcal C_\alpha(L,L')
  =(L^\dagger)^{\alpha/2}(L'^\alpha-L^\alpha)
                         (L^\dagger)^{\alpha/2}.
\end{equation}
It is zero on the common kernel. Its positive eigenvalues are
$q_i-1$, where the $q_i>1$ are the nontrivial relative eigenvalues on
$S$. A path edge is a bridge, so $b^TL^\dagger b=1/w$; its leverage
is exactly $\tau$.

\begin{theorem}[Sharp rank-two path trace bound]\label{thm:sharp}
Fix $\tau>0$ and set $K=\sqrt{1+\tau}$. Suppose a single edge increase
by the factor $1+\tau$ on a finite unit-mass weighted path has
$\rank\mathcal C_{1/2}=2$. Let $q_1\ge q_2>1$ be its two nontrivial
relative eigenvalues. Then $q_1q_2=K$ and
\begin{equation}\label{eq:sharp}
 q_1+q_2\ge T_*(\tau):=
 \frac{K+1}{\displaystyle\sqrt{1+
   \frac{(K-1)^2}{(K+1+2\sqrt{2K})^2}}}.
\end{equation}
The bound is attained by a four-vertex path with equal outer
conductances and middle-to-outer conductance ratio $1/K$, editing
the middle edge.
\end{theorem}

At $\tau=1$, the exact expression in \eqref{eq:sharp} is approximately
$2.4080334082788246$. In particular,
\begin{equation}\label{eq:q2numerical}
 q_2\le\frac{T_*(1)-\sqrt{T_*(1)^2-4\sqrt2}}2
       =1.0157544894662653\ldots.
\end{equation}
The decimals illustrate exact formulas; the exclusion below uses
rational inequalities.

\begin{theorem}[An explicit matrix spectrum excluded from all paths]
\label{thm:excluded}
At $\alpha=1/2$ and $\tau=1$, the relative spectrum
\begin{equation}\label{eq:target}
 (q_1,q_2)=\left(\frac54,\frac{4\sqrt2}{5}\right)
\end{equation}
is attained by a two-dimensional positive definite rank-one pair.
No single edge increase on any finite unit-mass weighted path attains
this spectrum with all other correction eigenvalues exactly zero.
\end{theorem}

Thus even a target with distinct positive correction eigenvalues
may require a limit in the path realization problem. The path length
and conductance ratios in Theorem~\ref{thm:excluded} are unrestricted.

\section{Rank and the complete classification of eligible paths}

\begin{lemma}[Rank equals cyclic dimension]\label{lem:rank}
For $A\succ0$, $u\ne0$, and $0<\alpha<1$,
\[
 \rank C_\alpha(A,u)
   =\dim\spanof\{u,Au,A^2u,\ldots\}.
\]
For \eqref{eq:graphcorrection}, this equals
$\dim\spanof\{b,Lb,L^2b,\ldots\}$.
\end{lemma}
\begin{proof}
The fractional resolvent representation and rank-one inverse identity
give
\begin{equation}\label{eq:resolvent}
 B^\alpha-A^\alpha
 =\frac{\sin(\pi\alpha)}\pi\int_0^\infty
 \frac{t^\alpha(A+tI)^{-1}uu^T(A+tI)^{-1}}
      {1+u^T(A+tI)^{-1}u}\,dt.
\end{equation}
The integral converges in norm and has positive semidefinite integrands.
A vector $x$ lies in its kernel precisely when
$u^T(A+tI)^{-1}x=0$ for almost every $t>0$. This analytic scalar
function must then vanish for all $t>0$. Expanding at infinity shows
that $u^TA^kx=0$ for every $k\ge0$. The converse follows by restricting
the resolvent to the invariant cyclic space
$\mathcal K=\spanof\{u,Au,\ldots\}$.

Hence the difference has kernel $\mathcal K^\perp$. Since $A$ is
symmetric, both $\mathcal K$ and $\mathcal K^\perp$ are invariant under
$A^{-\alpha/2}$. The congruence defining $C_\alpha$ preserves this
kernel and therefore its rank. Restriction to $S$ proves the graph
statement.
\end{proof}

\begin{lemma}[All rank-two edited paths]\label{lem:paths}
A finite unit-mass weighted path has cyclic dimension two at an
edited edge if and only if it is either
\begin{enumerate}
 \item a three-vertex path, with either edge edited; or
 \item a four-vertex path, with the middle edge edited and equal
       outer conductances.
\end{enumerate}
\end{lemma}
\begin{proof}
Let $r$ be the greatest distance from a path vertex to the set of
edited-edge endpoints. Along an arm of length $r$, the vectors
$b,Lb,\ldots,L^rb$ are linearly independent: at each new distance
$k$, $L^kb$ has a nonzero coordinate equal, up to sign, to the product
of the $k$ intervening positive conductances; all earlier powers
vanish there. Cyclic dimension two therefore requires $r\le1$.
There are at most two vertices on either side of the edited edge,
so the entire path has at most four vertices.

A two-vertex path has cyclic dimension one. For three vertices,
$Lb$ has a nonzero coordinate outside the support of $b$, so the
cyclic dimension is two. For four vertices the edited edge must be
the middle one. Write the conductances as $(a,m,c)$ and
$b=(0,1,-1,0)^T$. Then
\[
 Lb=(-a,a+2m,-c-2m,c)^T.
\]
If the cyclic dimension is two, $L^2b=sLb+pb$ for some real $s,p$.
The first and fourth coordinates give
\[
 s=2(a+m),\qquad s=2(c+m),
\]
so $a=c$. Conversely, if $a=c$, reflection preserves $L$ and its
two-dimensional antisymmetric subspace contains $b$. Again $Lb$ is
not a multiple of $b$, so the cyclic dimension is exactly two.
\end{proof}

\section{The two-dimensional calculation}

An orthogonal change of basis and a common positive scaling of both
endpoints preserve their relative eigenvalues. On the active space,
the two path classes in Lemma~\ref{lem:paths} reduce to
\begin{equation}\label{eq:canonical}
 A=\begin{pmatrix}
       1+\kappa x&-\sqrt\kappa x\\
       -\sqrt\kappa x&x
    \end{pmatrix},\qquad
 B=A+\tau e_1e_1^T,\qquad x>0,
\end{equation}
with $\kappa=1/3$ for three vertices and $\kappa=1$ for four vertices.
Indeed $\det A=x$ and $(A^{-1})_{11}=1$, so the leverage is $\tau$.

For clarity, the reductions can be checked directly. On three
vertices, let the edited conductance be $a$ and the other one $c$.
In a basis starting with the normalized edited-edge vector, the
baseline is
\[
 \begin{pmatrix}2a+c/2&-\sqrt3 c/2\\
                -\sqrt3 c/2&3c/2\end{pmatrix}.
\]
Dividing by $2a$ gives \eqref{eq:canonical} with
$x=3c/(4a)$ and $\kappa=1/3$. On four vertices, with equal outer
conductances $a$ and middle conductance $m$, the antisymmetric
baseline is
\[
 \begin{pmatrix}a&-a\\-a&a+2m\end{pmatrix}.
\]
Reorder its basis and divide by $2m$: then $x=a/(2m)$ and $\kappa=1$.
In both cases the scaled update is exactly $\tau e_1e_1^T$.

\begin{lemma}[Square-root trace formula]\label{lem:trace}
For \eqref{eq:canonical}, set $K=\sqrt{1+\tau}$,
$h=1+\kappa$, and $s=\sqrt x$. Then
\begin{align}
 \det Q_{1/2}&=K,\nonumber\\
 \tr Q_{1/2}&=(K+1)\frac{N}{\sqrt P},\label{eq:trace}\\
 N&=hs^2+(K+1)s+K,\nonumber\\
 P&=(hs^2+2s+1)(hs^2+2Ks+K^2)
    =N^2+\kappa(K-1)^2s^2.\nonumber
\end{align}
\end{lemma}
\begin{proof}
For any positive definite two-dimensional matrix $X$, Cayley--Hamilton
gives
\[
 X^{1/2}=\frac{X+\sqrt{\det X}\,I}
                  {\sqrt{\tr X+2\sqrt{\det X}}}.
\]
Here is a general trace identity that makes the substitution explicit.
For $B=A+uu^T$, write $t=\tr A$, $d=\sqrt{\det A}$,
$e=\norm u^2$, and $\tau=u^TA^{-1}u$. The two-dimensional inverse
formula gives
$u^TA^{-1/2}u=(e+\tau d)/\sqrt{t+2d}$. Thus
\begin{equation}\label{eq:generaltrace}
 \tr(A^{-1/2}B^{1/2})
 =\frac{(K+1)t+(K+1)^2d+e}
        {\sqrt{(t+2d)(t+e+2Kd)}}.
\end{equation}
The trace equals $\tr Q_{1/2}$ by cyclic invariance; no commutation
is assumed. In \eqref{eq:canonical}, $t=1+hs^2$, $d=s$, and $e=\tau$.
Substitution gives \eqref{eq:trace}. Multiplication of the two
quadratics proves the stated identity for $P$.
\end{proof}

Since
\[
 \frac Ns=hs+(K+1)+\frac Ks\ge K+1+2\sqrt{hK},
\]
formula \eqref{eq:trace} gives the exact family minimum
\begin{equation}\label{eq:familymin}
 \min_{x>0}\tr Q_{1/2}
 =\frac{K+1}{\displaystyle\sqrt{1+
   \frac{\kappa(K-1)^2}
        {(K+1+2\sqrt{(1+\kappa)K})^2}}},
 \qquad x=\frac K{1+\kappa}.
\end{equation}
To compare the two families, differentiate
\[
 \frac{\sqrt\kappa}{K+1+2\sqrt K\sqrt{1+\kappa}}.
\]
The derivative has the sign of
$K+1+2\sqrt K/\sqrt{1+\kappa}$, which is positive. Consequently
the minimum in \eqref{eq:familymin} decreases with $\kappa>0$.
The four-vertex family, $\kappa=1$, gives the global minimum.
Its minimizing $x=K/2$ is equivalent to $m/a=1/K$.
Together with Lemmas~\ref{lem:rank} and~\ref{lem:paths}, this proves
Theorem~\ref{thm:sharp}.

\section{The excluded pair and an exact matrix witness}

First set $K=\sqrt2$. For either admissible path family,
$0<\kappa\le1$ and
\[
 K+1+2\sqrt{(1+\kappa)K}>K+3.
\]
It follows from \eqref{eq:familymin} that every rank-two path satisfies
\begin{equation}\label{eq:rationalgap}
 \tr Q_{1/2}>
 \frac{(K+1)(K+3)}{\sqrt{(K+3)^2+(K-1)^2}}
 =\frac{5+4K}{\sqrt{14+4K}}>\frac{12}5.
\end{equation}
Squaring positive quantities, the last inequality reduces to
$424\sqrt2>591$, which holds because
$2\cdot424^2=359552>349281=591^2$.
In contrast, the target in \eqref{eq:target} satisfies
\[
 \frac54+\frac{4\sqrt2}5<\frac{12}5,
\]
as $\sqrt2<23/16$. This excludes it from every path with exactly
two positive correction eigenvalues.

To verify exact matrix feasibility without any general inverse-spectrum
theorem, permit arbitrary $\kappa>0$ in \eqref{eq:canonical}.
Write
\[
 K=\sqrt2,\quad c=K+1,\quad \beta=\sqrt K,
 \quad T=\frac54+\frac{4K}5,
\]
and define
\begin{equation}\label{eq:witness}
 v=\frac{c^2-2\beta T}{c(T-2\beta)},\qquad
 \kappa=v^2-1,\qquad x=\frac K{1+\kappa}.
\end{equation}
The unequal positive target entries have product $K$, so
$2\beta<T<c$. The denominator defining $v$ is positive, and
\[
 v-1=\frac{(c+2\beta)(c-T)}{c(T-2\beta)}>0.
\]
Hence $\kappa,x>0$. The matrix $A$ in \eqref{eq:canonical} is
positive definite, and $B=A+e_1e_1^T$ has leverage one.

At the minimizing $x$ in \eqref{eq:witness}, formula
\eqref{eq:familymin} simplifies to
\[
 \tr Q_{1/2}=c\frac{c+2\beta v}{cv+2\beta}=T.
\]
For example, this follows by expanding
$(c+2\beta v)^2+(v^2-1)(K-1)^2=(cv+2\beta)^2$.
The defining equation for $v$ then gives the desired trace exactly.
The determinant is $K$, so the two eigenvalues are exactly
$5/4$ and $4\sqrt2/5$. This proves Theorem~\ref{thm:excluded}.

\section{Scope, attribution, and verification}

The obstruction distinguishes exact realization from approximation.
Longer paths can have additional positive correction eigenvalues
tending to zero, and so need not have cyclic dimension two. Such
approximations are compatible with the path closure construction in
the companion repository manuscript
\texttt{relative-fractional-response.tex}. Exact zeros impose the
rank restriction proved here.

There is a second elementary restriction: the target pair already
uses the entire determinant budget $q_1q_2=\sqrt2$. Keeping these
two entries exact while adding positive correction eigenvalues at
the same leverage is impossible by \eqref{eq:budget}. This note
does not address exact realization on other graph classes or with
arbitrary positive vertex masses.

The fractional resolvent formula and rank-one inverse identity are
classical. So are cyclic subspaces, the two-dimensional square-root
formula, and the determinant lemma. Low-rank matrix-function and square-root updates
are established subjects~\cite{BCKS,SDA}; matrix functional calculus
is treated systematically in~\cite{Higham}. The candidate contribution
here is the complete rank-two path classification applied to this
relative metric, its sharp trace formula, and the explicit separation
from positive definite matrix realizability. No historical novelty
claim has been certified.

The source proof is preserved in
\texttt{exact-unitmass-path-exploration.txt}; a separate informal
review appears in \texttt{exact-path-rank-two-independent-review.txt}.
The diagnostic program \texttt{scripts/check\_rank\_two\_path\_gap.py}
checks the trace formulas on 920 cases and an independently specified
exact-radical matrix witness. Floating-point diagnostics corroborate
the formulas; the arguments above establish the result.

\begin{thebibliography}{9}
\bibitem{BCKS}
B. Beckermann, A. Cortinovis, D. Kressner, and M. Schweitzer.
Low-rank updates of matrix functions II: Rational Krylov methods.
\emph{SIAM J. Numer. Anal.} 59 (2021), 1325--1347.
\url{https://arxiv.org/abs/2008.11501}.
\bibitem{SDA}
S. Shmueli, P. Drineas, and H. Avron.
Low-rank updates of matrix square roots.
\emph{Numer. Linear Algebra Appl.}, e2528 (first published 2023).
\url{https://arxiv.org/abs/2201.13156}.
\bibitem{Higham}
N. J. Higham.
\emph{Functions of Matrices: Theory and Computation}.
SIAM, 2008.
\end{thebibliography}
\end{document}
